\begin{figure}[ht]
\centering
\begin{tikzpicture}[x=1cm, y=1cm, font=\small]

% Axes
\draw[->, thick] (0, 0) -- (7, 0) node[anchor=north east] {stated confidence};
\draw[->, thick] (0, 0) -- (0, 6) node[anchor=south east] {empirical hit-rate};

% Tick labels: 0, 0.5, 0.9, 0.95, 0.99
\foreach \x/\lbl in {0/0, 3/0.5, 5.4/0.9, 5.7/0.95, 5.94/0.99} {
  \draw (\x, -0.05) -- (\x, 0.05);
  \node[anchor=north, font=\scriptsize] at (\x, -0.05) {\lbl};
}
\foreach \y/\lbl in {0/0, 3/0.5, 5.4/0.9, 5.7/0.95, 5.94/0.99} {
  \draw (-0.05, \y) -- (0.05, \y);
  \node[anchor=east, font=\scriptsize] at (-0.05, \y) {\lbl};
}

% Perfect-calibration diagonal
\draw[thick, dashed, engblue] (0, 0) -- (6, 6) node[anchor=south west, font=\footnotesize] {perfect};

% Overconfident calibration (typical untrained engineer): below the line at high confidence
\draw[thick, engred]
  plot[smooth] coordinates
   {(0, 0.0) (1, 0.7) (2, 1.5) (3, 2.3) (4, 3.0)
    (5, 3.6) (5.4, 3.95) (5.7, 4.15) (5.94, 4.35)};
\node[anchor=west, engred, font=\footnotesize] at (5.94, 4.35) {overconfident};

% Underconfident: above the line at low confidence
\draw[thick, engblue]
  plot[smooth] coordinates
   {(0, 0.5) (1, 1.4) (2, 2.4) (3, 3.4) (4, 4.3) (5, 5.1) (5.4, 5.45) (5.7, 5.7) (5.94, 5.94)};
\node[anchor=east, engblue, font=\footnotesize] at (1.0, 1.5) {under-};
\node[anchor=east, engblue, font=\footnotesize] at (1.0, 1.0) {confident};

\end{tikzpicture}
\caption[Calibration plot for a population of confidence
intervals]{The reader who states 90\,\% confidence intervals on
each of fifty Fermi answers can plot, after the truth is in, the
fraction of intervals that captured truth at each stated
confidence. A well-calibrated estimator's curve lies on the
diagonal. Most untrained estimators sit below the diagonal at high
confidence (the overconfident curve), capturing truth only $70$ to
$80\,\%$ of the time when they claim $95\,\%$ confidence.
Calibration training, of the kind documented in the Good Judgment
Project \cite{text:ch09-tetlock-superforecasting-2015}, pulls the
curve toward the diagonal within a few months of disciplined
practice.}
\label{fig:vol01:ch09:calibration-plot}
\end{figure}
